\section{Continuous families of idempotents}

The statements of this section are finite-dimensional ingredients of the
continuity argument of~\cite[Section~II.F.2]{Schuch2011Phases}.

\begin{lemma}[Polynomial correction of an idempotent]
    \label{thm:spt_newton_idempotent}
    \lean{idempotent_newton_defect, idempotent_newton_sub,
        idempotent_newton_twice_sub_one, idempotent_newton_increment_bound,
        idempotent_newton_tendsto_of_small_defect}
    \leanok
    In a complete normed ring with $\|1\|\le1$, set
    $x_{n+1}=3x_n^2-2x_n^3$ and $d_n=x_n^2-x_n$.
    Then $d_{n+1}=d_n^2(4d_n-3)$ and
    $x_{n+1}-x_n=-(2x_n-1)d_n$.
    If $\|d_0\|\le1/16$, the sequence converges to an idempotent.
    More precisely,
    \begin{align}
        \|x_{n+1}-x_n\| & \le
        \|2x_0-1\|\,\|d_0\|\,2^{-n}.
        \notag
    \end{align}
\end{lemma}

\begin{proof}
    \leanok
    Polynomial expansion gives the identities, including
    $2x_{n+1}-1=(2x_n-1)(1-2d_n)$.
    Induction gives $\|d_n\|\le4^{-n}\|d_0\|$ and
    $\|2x_n-1\|\le2^n\|2x_0-1\|$. The increments are summable.
    Completeness gives the limit; continuity of multiplication makes it
    idempotent.
\end{proof}

\begin{lemma}[Continuous local idempotents]
    \label{thm:spt_continuous_nearby_idempotent}
    \lean{Matrix.exists_local_continuous_idempotent_in_adjoin}
    \leanok
    Let $X(t)$ be a continuous family of complex square matrices,
    with $X(t_0)$ idempotent. On an open neighbourhood of $t_0$ there is
    a continuous family of idempotents $P(t)$, with $P(t_0)=X(t_0)$,
    belonging pointwise to the unital algebra generated by $X(t)$.
    No Hermiticity assumption is required.
\end{lemma}

\begin{proof}
    \leanok
    \uses{thm:spt_newton_idempotent}
    Shrink the neighbourhood so that the idempotency defect is at most
    $1/16$ and $X$ is bounded. Apply the preceding correction in the
    complete normed algebra of bounded continuous matrix functions.
    Each iterate belongs to the generated algebra, which is closed
    by finite dimensionality. At $t_0$ every iterate equals $X(t_0)$.
\end{proof}

\begin{lemma}[Rank of an idempotent]
    \label{thm:spt_idempotent_rank}
    \lean{Matrix.trace_eq_rank_of_idempotent,
        Matrix.exists_open_constantRank_of_idempotent}
    \leanok
    An idempotent complex matrix has trace equal to its rank.
    Consequently the rank of a continuous family of idempotents is
    locally constant, including for non-Hermitian idempotents.
\end{lemma}

\begin{proof}
    \leanok
    The associated linear map is a projection onto its range.
    Its trace is the dimension of that range. Continuity of the trace
    and discreteness of the integer ranks give local constancy.
\end{proof}

\begin{lemma}[A fixed section of a constant-rank family]
    \label{thm:spt_trace_quotient_fixed_section}
    \lean{Matrix.exists_fixedSection_of_rank,
        Matrix.exists_open_fixedSection_range}
    \leanok
    Let $G(t)$ be a continuous matrix family of constant rank $r$.
    At any base point $t_0$ there is a coefficient section $F$ such
    that $G(t_0)F$ is injective. For this fixed $F$, there is an open
    neighbourhood $S$ of $t_0$ on which $G(t)F$ remains injective and
    has the same range as $G(t)$.
\end{lemma}

\begin{proof}
    \leanok
    Choose coordinates on the range of $G(t_0)$ and lift them through
    $G(t_0)$. The Hermitian Gram matrix of the resulting section is
    invertible. Invertibility persists on an open neighbourhood.
    The section range is contained in the full range and both have
    dimension $r$, so they coincide there.
\end{proof}

\begin{lemma}[Gram left inverses]
    \label{thm:spt_gram_left_inverse}
    \lean{Matrix.gramLeftInverse_mul, Matrix.continuous_gramLeftInverse}
    \leanok
    If a complex rectangular matrix $C$ has independent columns, then
    $H=(C^\dagger C)^{-1}C^\dagger$ satisfies $HC=I$.
    For a continuous family with independent columns, $H$ is continuous.
\end{lemma}

\begin{proof}
    \leanok
    The Gram matrix is positive definite and hence invertible.
    Matrix inversion is continuous on the invertible matrices.
\end{proof}

\begin{lemma}[Continuous frames for a matrix range]
    \label{thm:spt_continuous_idempotent_frame}
    \lean{Matrix.exists_open_continuous_rangeFrame_of_constantRank,
        Matrix.exists_local_continuous_rangeFrame_of_idempotent}
    \leanok
    A continuous matrix family $P(t)$ of constant rank $D$ admits a local
    continuous frame $C(t)=P(t)F$ for its range, with a fixed matrix $F$.
    The matrices $H(t)=(C(t)^\dagger C(t))^{-1}C(t)^\dagger$ are continuous
    and satisfy $H(t)C(t)=I_D$.
    For a continuous family of idempotents it suffices that $P(t_0)$ have
    rank $D$; in addition $P(t)C(t)=C(t)$.
\end{lemma}

\begin{proof}
    \leanok
    \uses{thm:spt_trace_quotient_fixed_section, thm:spt_gram_left_inverse,
        thm:spt_idempotent_rank}
    Choose a section at the base parameter. Injectivity persists nearby.
    Equality of dimensions identifies its range with the matrix range.
    Apply the Gram left inverse and the local constancy of the rank.
\end{proof}

\begin{lemma}[Coordinates of a bilinear map]
    \label{thm:spt_bilinear_coordinates}
    \lean{Matrix.toLinearMap₂'_apply_mulVec_prod}
    \leanok
    Let $M$ have columns indexed by pairs $(a,b)$. The associated bilinear
    map satisfies $\mu(x,y)=M(x\otimes y)$, where
    $(x\otimes y)_{(a,b)}=x_a y_b$.
\end{lemma}

\begin{proof}
    \leanok
    Expand both sides in the coordinate basis.
\end{proof}

\begin{lemma}[Dimension of a primitive left ideal]
    \label{thm:spt_primitive_left_ideal_dimension}
    \lean{Matrix.finrank_range_mulRight_single,
        Matrix.finrank_range_rightMul_of_matrixAlgebra}
    \leanok
    In $\MN{D}$ the left ideal $\MN{D}E_{ii}$ has
    dimension $D$. The same holds under any algebra
    isomorphism.
\end{lemma}

\begin{proof}
    \leanok
    The ideal consists of matrices supported in column $i$.
    That column is an arbitrary vector in $\C^D$.
\end{proof}
